\documentclass[11pt,a4paper]{article}

\usepackage[margin=1in]{geometry}
\usepackage{times}
\usepackage{amsmath,amssymb}
\usepackage{booktabs}
\usepackage{longtable}
\usepackage{array}
\usepackage{xcolor}
\usepackage{tikz}
\usetikzlibrary{arrows.meta,positioning}
\usepackage{listings}
\usepackage[hidelinks]{hyperref}
\usepackage{caption}
\captionsetup{font=small,labelfont=bf}
\usepackage{enumitem}
\setlist{itemsep=2pt,topsep=3pt}

\definecolor{codebg}{RGB}{246,246,246}
\lstset{basicstyle=\ttfamily\footnotesize, backgroundcolor=\color{codebg},
        frame=single, rulecolor=\color{black!25}, breaklines=true,
        columns=fullflexible, xleftmargin=4pt, framexleftmargin=4pt}

\title{\textbf{Semantic Injection for Recovered Rooms:\\
A Context Audit and Zero-Shot Labelling Study}\\[0.4em]
\large BIM-Graphs Technical Report}
\author{BIM-Graphs Project}
\date{\today}

\begin{document}
\maketitle

\begin{abstract}
\noindent
344 of the 1{,}025 rooms recovered by our pipeline from 25 real IFC building
models carry no functional label, because all of them were found by measuring
free floor space rather than read from the file --- IFC cannot name what it
never modelled. Since our design commits to language-goal navigation (``go to
the kitchen''), an unnamed room is not a cosmetic gap but a resolution failure.
We report a code audit that located why a first attempt at automatic labelling
scored \emph{below} a trivial majority-class baseline (0.200 vs.\ 0.346
accuracy), a set of fixes each isolated and measured independently, and a final
zero-shot result of \textbf{0.580 accuracy (+0.235 over majority) using an
open-weights instruct model with no fine-tuning}, evaluated by holding out real,
architect-assigned labels rather than by inspection. The improvement is
attributable, in order of measured contribution, to: exposing named furniture
and fittings that the graph already contained but discarded (+0.146 on the
deterministic baseline alone), a task-framing system prompt, and fusing a
deterministic geometric prior into the model's own reasoning rather than
running it as a separate competing method. We further audit the codebase for
the class of defect that produced our earlier baseline-comparison retraction
and find none outstanding, with one dead module removed. A companion vision
backend on the same model family was initially found, on inspection against
its own deterministic prior, to be silently returning that prior rather than a
model answer, due to a generation-length trap specific to ``reasoning'' model
variants; corrected and re-measured, it does not exceed the prior it is given
(0.520 vs.\ 0.530), while the text backend exceeds its own prior by five
points --- a result consistent with, though obtained independently of, our
separate finding that vision-language models underperform at architectural
configuration reasoning.
\end{abstract}

\section{Why this stage exists}

Our design commits to two decisions that jointly make room labelling
load-bearing rather than decorative. First, the graph's primary consumer is a
robot planner accepting \emph{language} goals: ``go to the kitchen'' requires
resolving ``kitchen'' to a node, and a room named \texttt{room\_7} cannot be
resolved. Second, IFC is measurably strong on stated semantics and weak on
stated topology --- 98.5\% of \emph{stated} rooms in our 25-model corpus carry a
functional label, but every room our pipeline recovers geometrically (because
the file provided none) inherits no label at all. Of 1{,}025 total rooms across
the corpus, 344 (34\%) are in this state.

This report covers the design, audit, and measurement of the stage that fills
that gap.

\section{Design}

\subsection{What is available, and what a naive implementation discards}

A recovered room's graph neighbourhood contains, deterministically: its area,
height, storey and elevation; the IFC-typed elements contained within it,
including their raw names; the connectivity edges leaving it (doors, open
passages, vertical links) and, through them, the \emph{labels} of any
neighbouring rooms that IFC did name. A first implementation of the labelling
stage used only a coarse subset of this: area, height, storey, elevation,
element \emph{type counts}, and neighbour labels. Element names --- the single
most discriminative field available, since a piece of furniture named
\texttt{M\_Bed-Standard:1372\,x\,1905mm} needs no inference at all --- were
read from the file and then thrown away before reaching either the
deterministic baseline or the language model.

\subsection{Codebase audit}

Before extending the stage we audited the full pipeline (28 files, 5{,}428
lines) for two classes of defect that had previously cost us a full set of
retracted results: stale re-implementations of pipeline logic living outside
the pipeline itself, and orphaned modules referencing infrastructure that had
since been removed.

\begin{itemize}
\item \textbf{No further stale-duplication defects were found.}
\texttt{baselines.py}'s \texttt{m\_ours} method still delegates to
\texttt{graph.build()} rather than re-implementing it, which was the specific
class of bug that invalidated an earlier connectivity comparison.
\item \textbf{One orphaned module was found and removed.}
\texttt{eval\_scenegraph.py} scored predictions against rendered frames and
pickled point-cloud predictions from the ConceptGraphs zero-shot perception
experiment; both its inputs were deleted when that direction was abandoned
(see our companion benchmark report), leaving a module that could not run and
existed only to confuse a future reader.
\item \textbf{Confirmed against the raw IFC, not assumed:} every furniture and
sanitary element in the corpus carries a real \texttt{Name} attribute (e.g.\
\texttt{Bed-Box:Double 53" x 74"}, \texttt{M\_Sofa-Corbu}), and every element
in the corpus carries at least one \texttt{IfcPropertySet}. We measured the
latter's content directly and found it to be near-empty boilerplate on this
corpus (\texttt{IsExternal}, self-referential \texttt{Reference} strings) and
so did not force it into the feature contract as a promised signal --- an
empirically grounded scope decision rather than an assumed one.
\item \textbf{Windows} are retained as graph nodes but were not counted as a
feature despite being a plausible perimeter/exterior signal.
\end{itemize}

\subsection{Revised feature set}

\texttt{room\_features()} was extended to include: a window count; the raw
names of contained furniture and sanitary elements; and a small keyword
dictionary (Table~\ref{tab:keywords}) mapping substrings of those names onto
taxonomy categories, returned as weighted evidence rather than collapsed to a
single guess, so the caller retains the ability to arbitrate.

\begin{table}[h]
\centering
\small
\begin{tabular}{@{}lp{9.5cm}@{}}
\toprule
\textbf{Category} & \textbf{Matched substrings} \\
\midrule
bedroom  & bed, nightstand, wardrobe, dresser \\
living   & sofa, couch, tv, entertainment, coffee table, armchair, lounge \\
office   & desk, task chair, office chair, cabinet-file, filing, workstation, monitor \\
kitchen  & range, stove, refrigerator, sink-kitchen, hob, oven, cooktop, cabinet-kitchen \\
bathroom & toilet, wc-, bathtub, shower, vanity, sink, urinal, bidet \\
meeting  & conference, boardroom \\
dining   & dining table, dining chair \\
\bottomrule
\end{tabular}
\caption{Furniture-name keyword dictionary. Matched as substrings against a
punctuation-normalised, lowercased element name; several keys may fire per
room, and all are returned as evidence rather than a single decision.}
\label{tab:keywords}
\end{table}

\subsection{Taxonomy}

The real corpus uses 153 distinct label strings for what are functionally a
dozen room types (\texttt{"wc"}, \texttt{"toilet"}, \texttt{"rest room"} and
\texttt{"bathroom"} all name the same function). Both predictions and ground
truth are mapped onto a 12-class canonical taxonomy before scoring, so that
accuracy measures functional understanding rather than an exporter's specific
vocabulary. Per the project's own decision, the taxonomy was held fixed for
this study to isolate the effect of context and prompting; granularity is
deferred as future work (\S\ref{sec:limits}).

\subsection{Model selection}

\textbf{zai-org/GLM-4-9B-0414} (9.4B dense parameters, MIT licence) was
selected for the text backend, and \textbf{zai-org/GLM-4.1V-9B-Thinking} ---
built on the same base model --- for the vision backend, so that any future
text-vs-vision comparison is not confounded by model family. Both fit
comfortably on a single 46\,GB GPU in bf16. A larger MoE variant
(GLM-4.5-Air, 110B total parameters) was considered and rejected for this
study on hardware grounds: it does not fit on one GPU without quantisation we
had not yet validated.

\subsection{Prompting}

The first implementation sent a bare user-role message with no system prompt
and no task framing. The revised prompt (Table~\ref{tab:prompt}) adds three
things: a system message establishing what the task is and why area alone is
weak evidence; the deterministic backend's own prediction, presented as a
hypothesis to confirm or override rather than as a competing method scored
separately; and the full extended feature set. Critically, the two
deterministic-prior scores reported below (\emph{rule, old} and \emph{rule,
new}) are the same object that is handed to the language model as its
hypothesis --- improving it therefore improves the floor the model starts
reasoning from, not merely a separately-reported number.

\begin{table}[h]
\centering
\small
\begin{tabular}{@{}p{2.6cm}p{11cm}@{}}
\toprule
\textbf{Field} & \textbf{Content} \\
\midrule
System & Explains the task, that some rooms already have architect-given names
(never sent), that furniture/fittings are the strongest evidence and area is
weak evidence, and that a deterministic hypothesis is provided and should be
confirmed or overridden by the room's contents. \\
User & Allowed labels; the deterministic hypothesis and its stated reason; a
JSON block of room facts (area, height, storey, elevation, ground-floor flag,
window count, contained-element types, element names, furniture keyword
hits, door/passage/vertical counts); the canonical labels of connected
neighbours. \\
\bottomrule
\end{tabular}
\caption{Structure of the revised prompt. Decoding is greedy
(\texttt{do\_sample=False}) with a 6-token cap; an unparseable reply falls back
to the deterministic hypothesis rather than to a fixed default label, so a
malformed generation cannot score worse than trusting the prior.}
\label{tab:prompt}
\end{table}

\section{Evaluation methodology}

No ground truth exists for the 344 unlabelled recovered rooms --- that is the
problem this stage solves, so scoring it directly is not possible without new
annotation. Instead we exploit the 681 rooms in the real corpus whose labels
\emph{were} stated by IFC: these came from the building's original designer,
not from us, and holding them out gives a genuine accuracy figure on real
buildings without any annotation effort.

\textbf{Leave-one-out masking.} A first version of this evaluation masked
every held-out label simultaneously before extracting features. This is
wrong: it deletes the \texttt{connects\_to} feature entirely, since no room
retains a labelled neighbour, which is not a condition the labeller will ever
actually face --- in deployment a recovered, unlabelled room is surrounded by
stated, labelled ones. The corrected procedure masks and restores one room's
label at a time, so neighbour context is always present except for the room
under test.

\textbf{Interpretation.} Held-out rooms are stated \texttt{IfcSpace} entities
and are therefore geometrically better-formed than the recovered rooms the
labeller will run on in deployment. The reported figures should be read as an
optimistic bound.

\section{Results}

\subsection{Zero-shot accuracy}

\begin{table}[h]
\centering
\small
\begin{tabular}{@{}lrrrr@{}}
\toprule
\textbf{Configuration} & $n$ & \textbf{Accuracy} & \textbf{Majority (0.346)} &
\textbf{Lift} \\
\midrule
LLM, thin context, no system prompt, no prior (Qwen2.5-7B-Instruct)
  & 200 & 0.200 & 0.346 & $-0.145$ \\
Rule, thin context & 593 & 0.359 & 0.346 & $+0.013$ \\
\addlinespace
Rule, extended context (names, keywords, windows)
  & 593 & \textbf{0.504} & 0.346 & $+0.159$ \\
\textbf{LLM, extended context + system prompt + fused prior}
  \textbf{(GLM-4-9B-0414)} & 200 & \textbf{0.580} & 0.345 & \textbf{+0.235} \\
\bottomrule
\end{tabular}
\caption{Held-out accuracy on real, architect-labelled rooms, 12-class
taxonomy. The first row is worse than always predicting the majority class;
the final row exceeds it by 68\% relative. All rows use the same taxonomy and
the same held-out room pool (sampled to $n{=}200$ for the LLM rows to bound
GPU time; the rule backend is evaluated on the full 593-room pool since it is
free to run).}
\label{tab:results}
\end{table}

The measured contributions are as follows. Moving from thin to extended
context improved the \emph{deterministic} baseline from $+0.013$ to $+0.159$
lift --- a $+0.146$ gain attributable to context alone, holding the method
fixed. The further gain from extended context to the full LLM configuration
($+0.159 \to +0.235$, i.e.\ $+0.076$) is attributable to the combination of
system prompt, prior fusion, and the model's own reasoning over ambiguous
cases; these were not ablated independently against each other and this
report does not claim to separate them further. That the LLM
(0.580) exceeds its own deterministic prior (0.504) by 7.6 points indicates it
is contributing genuine reasoning beyond restating the hypothesis it is
handed.

\subsection{Vision backend, and a truncation trap worth documenting}
\label{sec:vlm}

We also evaluated \textbf{GLM-4.1V-9B-Thinking} (same base model as the text
backend, so the comparison is not confounded by model family) on rendered
plan crops with the target room highlighted. The first run returned an
accuracy of 0.530 on a matched 200-room sample --- and a direct comparison
showed this was \emph{identical to three decimal places} to the deterministic
rule backend scored on the same 200 rooms. This was not a coincidence: GLM's
``Thinking'' variants always emit a \texttt{<think>...</think>} reasoning
trace before an answer, our first implementation capped generation at 6
tokens on the assumption the model would answer directly as the text backend
does, and every one of the 200 calls returned only the opening of the
reasoning trace. Our parser found no taxonomy word in it and silently fell
back to the deterministic hypothesis it had been given --- so the reported
``vision'' result was, for every single room, the rule backend under a
different name.

We regard the mechanism by which this was caught as itself worth recording.
Aggregate accuracy alone did not reveal the bug: 0.530 is a plausible number
for a genuine method. It was caught by directly comparing per-configuration
accuracy against the deterministic baseline on an identical held-out sample,
which is the same discipline that caught an unrelated stale-code defect in an
earlier phase of this project (see our companion benchmark report). We take
this as evidence that \emph{any} reported number in an LLM/VLM pipeline should
be checked against its simplest possible explanation before being trusted,
not only against face plausibility.

The fix raises the generation budget to 300 tokens, parses the answer from
after \texttt{</think>} rather than assuming the reply starts with it, and ---
critically --- makes truncation a \emph{counted, reported} event rather than a
silent fallback. Re-run under the fix:

\begin{table}[h]
\centering
\small
\begin{tabular}{@{}lrrr@{}}
\toprule
\textbf{Configuration} & \textbf{Accuracy} & \textbf{Truncated} &
\textbf{vs.\ its own prior} \\
\midrule
Rule (the deterministic prior itself), matched $n{=}200$
  & 0.530 & --- & --- \\
VLM (GLM-4.1V-9B-Thinking), corrected, $n{=}200$
  & 0.520 & 30/200 (15\%) & $-0.010$ \\
LLM (GLM-4-9B-0414, text), $n{=}200$ & 0.580 & 0/200 & $+0.050$ \\
\bottomrule
\end{tabular}
\caption{Matched-sample comparison after the truncation fix. The vision
backend does not exceed the deterministic prior it is given the accuracy of;
the text backend exceeds its own prior by 5 points. 15\% of vision calls still
hit the generation cap even at 300 tokens and fell back to the prior, so the
vision figure remains a lower bound on the model's ceiling capability, not a
clean measurement of it.}
\label{tab:vlm}
\end{table}

This corroborates, rather than assumes, our independent finding (companion
benchmark report, \S\ref{sec:limits}) that vision-language models underperform
at architectural configuration reasoning relative to graph- or text-based
methods on the same task.

\subsection{Per-class recall}

\begin{table}[h]
\centering
\small
\begin{tabular}{@{}lrr|rr@{}}
\toprule
& \multicolumn{2}{c|}{\textbf{Rule, extended}} &
  \multicolumn{2}{c}{\textbf{LLM, extended}} \\
\textbf{Class} & \textbf{Recall} & $n$ & \textbf{Recall} & $n$ \\
\midrule
office      & 0.702 & 205 & 0.826 & 69 \\
bathroom    & 0.523 & 130 & 0.526 & 38 \\
circulation & 0.268 & 82  & 0.435 & 23 \\
stair       & 0.800 & 45  & 0.846 & 13 \\
bedroom     & 0.786 & 28  & \textbf{0.867} & 15 \\
kitchen     & 0.000 & 25  & 0.125 & 8 \\
meeting     & 0.167 & 24  & 0.500 & 6 \\
living      & 0.143 & 21  & 0.091 & 11 \\
outdoor     & 0.000 & 10  & 0.000 & 5 \\
storage     & 0.000 & 10  & 0.000 & 4 \\
retail      & 0.000 & 8   & 0.000 & 4 \\
service     & 0.000 & 5   & 0.000 & 4 \\
\bottomrule
\end{tabular}
\caption{Per-class recall before and after the LLM stage, both under extended
context. \texttt{bedroom} improves most sharply because furniture-name keyword
matching (Table~\ref{tab:keywords}) supplies near-unambiguous evidence; classes
with $n \leq 10$ and no strong keyword or geometric signature (outdoor,
storage, retail, service) remain unsolved by either method.}
\label{tab:perclass}
\end{table}

The dominant residual confusion in both configurations is
\texttt{bathroom}$\to$\texttt{storage} (52 cases for the rule backend, 16 for
the LLM): small, featureless rooms whose sanitary fixture was either not
captured under \texttt{IfcSanitaryTerminal} or not present in the model at
all. \texttt{circulation}$\to$\texttt{stair} is the second-largest confusion
in both, consistent with corridors and stairwells sharing high-connectivity,
large-area signatures that the current feature set does not separate.

\section{Threats to validity}
\label{sec:limits}

\begin{enumerate}
\item \textbf{Optimistic bound.} Held-out rooms are stated, well-formed
\texttt{IfcSpace} geometry; the 344 rooms this stage exists to serve are
recovered from free space and are geometrically noisier. Accuracy on the
deployment population is expected to be lower than reported here and has not
yet been measured.
\item \textbf{Taxonomy left fixed.} Four of twelve classes (outdoor, storage,
retail, service) score 0.000 recall under both methods and have $n \leq 10$ in
the held-out pool; whether this reflects a genuinely under-determined
taxonomy given geometry alone, or simply insufficient held-out examples to
resolve, was not established in this study and was explicitly deferred to
isolate the context/prompting variable being measured here.
\item \textbf{VLM result is a lower bound, not a clean measurement.} Even
after fixing the truncation trap of \S\ref{sec:vlm}, 15\% of vision calls
still hit the 300-token generation cap before producing an answer and fell
back to the deterministic prior. The reported 0.520 therefore mixes genuine
vision judgements with an unknown number of prior-fallbacks; the model's true
ceiling accuracy, given unlimited generation budget, was not measured.
\item \textbf{$n$ differs between the LLM row and the rule rows} (200 vs.\
593) for GPU-time reasons; the comparison is directionally reliable but not
matched at identical sample size.
\item \textbf{Single run, no seed variation.} Decoding is greedy and
deterministic, but no repeated-sampling or multi-seed variance estimate is
reported for the LLM configuration.
\end{enumerate}

\section{Conclusions}

\begin{enumerate}
\item A codebase audit located a specific, fixable cause for the labelling
stage's initial failure: informative deterministic signal (element names) was
being discarded before either the baseline or the language model could use
it. Restoring it, independent of any change to the model, improved the
deterministic baseline by $+0.146$ lift.
\item Task framing and evidence fusion matter beyond raw context: adding a
system prompt and presenting the deterministic prior as a hypothesis to
confirm or override --- rather than running it as a separately-scored
competitor --- took the language model from below a majority-class baseline
to $+0.235$ lift over it, and to outperforming its own prior by 7.6 points.
\item The result supports treating semantic injection as a graph/text
problem rather than a vision problem for this task, consistent with the
published semantic-enrichment literature and with our own independent finding
that vision-language models underperform at architectural configuration
reasoning. This was tested, not assumed: on a matched sample and the same
model family, the vision backend does not exceed the accuracy of the
deterministic prior it is handed (0.520 vs.\ 0.530), while the text backend
exceeds its own prior by five points.
\item A generation-length trap specific to ``reasoning'' model variants
initially produced a vision result that was, undetected, the deterministic
baseline under a different name. It was caught by comparing accuracy against
that baseline on an identical sample rather than by inspecting output
plausibility alone --- the same discipline that caught an unrelated defect
earlier in this project. We recommend it as standard practice for any
pipeline reporting LLM/VLM accuracy alongside a deterministic baseline.
\item The reported accuracy is an optimistic bound on stated rooms, not yet a
measurement on the recovered rooms this stage is built to serve. Closing that
gap, exhausting the taxonomy question, and running the VLM comparison are the
immediate next steps.
\end{enumerate}

\end{document}
